\chapter{Conclusion}
\label{ch:conclusion}

This thesis presented three computational contributions for modelling oral biofilm
and microbiome dynamics, unified by the Extended Hamilton Principle.

\textbf{Chapter~\ref{ch:heine}} demonstrated that GPU-accelerated TMCMC enables
tractable Bayesian inference of a 15-dimensional interaction matrix from
limited in-vitro time-course data, achieving ${\sim}200\times$ speedup and
correctly recovering the qualitative distinction between commensal and dysbiotic
community dynamics in Heine et al.'s four experimental conditions.

\textbf{Chapter~\ref{ch:dieckow}} showed that three-layer metabolic sign priors
— constructed from experimental eHOMD data, AGORA2 pFBA, and MICOM community
FBA — consistently reduce out-of-sample prediction error in leave-one-patient-out
cross-validation (LOO-RMSE $= 0.0490$, Bray-Curtis $= 0.147$, 47.7\% below
the persistence null model), with external validation on 127 independent
peri-implant samples confirming ecological generalisation.

\textbf{Chapter~\ref{ch:integration}} carried the same interactions into space.
A reaction--diffusion extension of the Hamilton model was formulated, its
transport operators verified to be second order in the diffusion stencil and
mass-conserving only in conservative finite-volume form (the diagonal scheme's
ghost-node boundary is not), and confronted with the three-dimensional structure
extracted from $84$ FISH fields of view. The imaging delivered the spatial
work's principal empirical result: dysbiosis is a laterally homogeneous, vertically
stratified mat with \textit{P.~gingivalis} sequestered at depth beneath
\textit{F.~nucleatum} ($3$-D Manders overlap collapsing to $0.2$--$0.3$ while the
lateral projection reads near unity), whereas commensal health is patchy and
microcolonial ($p=2.2\times10^{-4}$ for the homogenisation contrast). The continuum
model reproduces the dysbiotic homogenisation but not the commensal patchiness,
sharply delimiting what a reaction--diffusion field can and cannot explain.
A finite-element realisation then closed the chain from composition to mechanics:
embedding the depth composition as the growth driver of a continuum model --- with the
coupling verified by patch, energy, manufactured-solution and consistent-tangent checks
--- showed that the dysbiotic structure concentrates a far higher residual stress at the
implant interface and, at identical mechanical parameters, delaminates the biofilm
$3.3\times$ further than the commensal control, a modulus-agnostic mechanical route from
compositional dysbiosis to implant detachment.

Together, these contributions advance the goal of \emph{physics-informed,
data-driven microbial community modelling}: by anchoring statistical inference
in thermodynamic variational principles and metabolic network constraints,
they provide biologically interpretable and computationally scalable methods
applicable to clinical settings where data are inherently sparse.

\section{Limitations}
\label{sec:concl_limitations}

Four limitations qualify the conclusions. First, both datasets are small:
$N=10$ patients with three weekly time points in vivo, and five species under four
conditions in vitro. The in-vivo metabolic-prior improvement over the
experimental-only prior is therefore only marginally significant ($p\approx0.07$)
and should be read as a consistency result rather than a powered effect. Second,
the Hamilton interaction matrix is constrained symmetric ($A_{ij}=A_{ji}$) by the
variational derivation; genuinely asymmetric ecological interactions --- where the
unconstrained gLV gains its small predictive edge --- cannot be represented, and
the metabolic sign prior usefully constrains interaction direction but not
magnitude (FBA flux units are not ecological interaction units). Third, the FISH imaging protocol is destructive (fixation kills the biofilm), so
only $T=5$ discrete timepoints are available; the temporal derivative
$\partial_t\phi_i$ is too coarsely sampled to identify species diffusivities $D_i$
quantitatively, and the spatial PDE is therefore validated structurally rather than
with measured transport coefficients.
Fourth, the in-vitro commensal and dysbiotic consortia employ distinct
\textit{P.~gingivalis} strains --- the avirulent ATCC\,20709 under commensal
conditions and the virulent W83 under dysbiotic conditions --- so that inferred
interaction differences partly reflect strain-specific virulence factors
(gingipain and fimbrial gene complement) rather than purely environmental
selection; disentangling the two effects would require cross-strain experiments
outside the present dataset.

\section{Outlook}
\label{sec:concl_outlook}

Several extensions follow directly from the three results chapters, ordered by
proximity to the existing framework.

\paragraph{Near term: robustness and unified inference.}
The sign enrichment result ($\mathrm{SA}=78.6\%$, $p=0.0004$) was obtained
with a warm-started no-prior LOO; a neutral-initialised repeat from
$\mathbf{A}=\mathbf{0}$ confirms the enrichment is data-driven
($\mathrm{SA}=71.4\%$, $p=0.027$), ruling out warm-start as a confound. Once confirmed, the two inference frameworks can be
unified: TMCMC posteriors regularised by the metabolic sign prior would give
full credible intervals on in-vivo interaction strengths, not merely point
estimates, even on sparse clinical data. The cross-feeding enrichment being
confined to positive interactions also motivates a targeted extension: adding an
asymmetric competition channel to the Hamilton model (relaxing $A_{ij}=A_{ji}$
for pairs where the metabolic evidence is directional) would test whether the
competition direction recovers significance that the symmetric constraint
currently suppresses.

\paragraph{Medium term: spatial and metabolic depth.}
The falsifiable in-vitro prediction --- omitting \textit{F.~nucleatum} should
collapse the dysbiotic \textit{Pg} surge --- is a direct co-culture test of the
inferred network mechanism. On the spatial side, the PINN-fitted diffusivities
($D_{Fn}$, $D_{Pg}$, advection $u$) are the first quantitative estimates from
FISH $z$-profiles; extending the fit to a full time-series (rather than
steady-state slices) and coupling an oxygen-diffusion PDE to the growth rates
$b_i$ via Monod kinetics would test whether the low $D_{Pg}$ reflects anaerobic
niche selection or genuine transport limitation. The observed commensal patchiness,
which the reaction--diffusion model cannot reproduce, motivates coupling the
continuum field to a stochastic surface-attachment process.

\paragraph{Longer term: clinical translation.}
The guild-level dysbiosis ordering validated on the Joshi and Szafrański cohorts
suggests a route to a clinically deployable dysbiosis index: a low-dimensional
projection of the 10-guild composition onto the commensal--dysbiotic axis,
calibrated against the inferred attractor states and interpretable in terms of
the Veillonella cross-feeding backbone. Crucially, such an index is interpretable only
as a \emph{relative ordering}, not on an absolute scale: on the longitudinal
periodontitis cohort every patient scores a class-level $\mathrm{GDI}<0$ (the abundant
commensal classes \textit{Streptococcus}, \textit{Veillonella} and \textit{Actinomyces}
numerically dominate even in disease), so a fixed $\mathrm{GDI}=0$ cut would mis-classify
the entire diseased cohort as stable; the predicted bone-loss \emph{ranking} is
nonetheless invariant to the threshold choice (Spearman $\rho=0.97$ between a
cohort-relative and an absolute cut), so a deployable index must be anchored to a
healthy-cohort reference rather than to an absolute zero. Patient-specific MAGs (metagenome-assembled
genomes) reconstructed from deeper sequencing would replace reference-strain AGORA2
models with personalised sign priors, tightening the credible intervals on $A_{ij}$
without requiring additional time points. In each direction the same principle
recurs: anchoring data-driven inference in mechanistic --- thermodynamic and
metabolic --- structure is what makes scarce microbiome data yield testable,
interpretable models.
